
\subsubsection{布鲁金斯学会：多维AI竞赛与中国的多赛道策略}

布鲁金斯学会在2026年4月发布的两份研究报告为本文的分析框架提供了重要支撑。《Competing AI Strategies for the US and China》提出，美中AI竞争已从单一维度的模型性能比拼演变为跨越算力、模型能力、行业应用、系统集成和全球部署的五维竞赛。在这一多维框架下，中美在\emph{不同维度上各有优势}：美国在前沿模型能力上保持领先，中国在应用扩散速度和开源生态渗透上大幅领先。

更具洞察力的是《China is running multiple AI races》一文的核心判断：中国AI开发者并非在简单追赶GPT-5——他们正在跑\emph{不同的赛道}。在产业约束（芯片出口管制）和北京政策焦点（AI赋能实体经济）的双重驱动下，中国AI产业形成了以开源生态、工业应用和成本优势为核心的差异化竞争力。这与本文光谱扩散框架的核心逻辑完全一致：降低获取门槛、加速光谱爬升，比追求单一模型的极致性能具有更大的社会影响力。

Brookings研究员Kyle Chan在2026年7月接受CNBC采访时进一步确认了这一趋势：''中国AI模型对美国公司特别有吸引力，因为AI成本正在飙升。过去美国公司优先考虑AI采用不论模型，现在他们越来越注重成本。他估计中国前沿开源模型目前落后美国顶级模型''六到九个月''——这一差距在成本优势面前正在快速缩小。

\subsubsection{兰德公司：开源作为软实力与AI竞争光谱}

兰德公司在《Open Models, Soft Power, and the Spectrum of U.S.-China Artificial Intelligence Competition》中提出了一个与本文高度共振的分析框架。报告的核心贡献有二：

第一，将开源AI定位为\textbf{软实力工具}。中国通过开源模型（DeepSeek、GLM、Qwen等）扩大对全球南方国家的影响力——这些国家无法负担OpenAI或Anthropic的高昂API费用，但可以通过部署中国开源模型获得AI能力。这是认知民主化在国际体系层面的具体体现：开源使AI能力从少数富裕国家的特权变为全球公共品。

第二，中美AI竞争不是一个二元对立，而是一个\textbf{光谱}——从完全开源到完全闭源的连续分布。中国位于光谱的开源端，美国在拜登时期向闭源端移动（受安全行政令推动），在特朗普第二任期又向开源端回摆。这一政策摇摆本身——兰德报告指出——可能损害美国在开源生态中的信誉和领导力。

兰德的政策建议对本文具有直接参考价值：\textbf{美国若过度限制开源，可能将全球AI生态拱手让给中国。}

\subsubsection{高盛AI采用率追踪：光谱分布的实证锚点}

高盛于2026年4月发布的AI Adoption Tracker为本文的光谱分布假说提供了最接近系统实证数据的支撑。基于美国人口普查局的商业趋势与展望调查（BTOS），高盛报告了以下关键数据：

\begin{center}
\begin{tabular}{c|c}
\hline
指标 & 数据 \\
\hline
美国企业整体AI采用率 & 19.8\%（预期6个月后22.3\%） \\
大型企业（$>）采用率 & 35.3\% \\
小型企业（20-49人）增速 & +2.1个百分点至21.5\% \\
信息技术/网络服务业采用率 & 60\%（行业最高） \\
AI每日节省工时 & 40--60分钟（OpenAI企业数据） \\
学术研究：生产力提升 & 平均23\% \\
企业案例：效率提升 & 约33\% \\
\hline
\end{tabular}
\end{center}

这些数据对本文的理论框架提供了三重验证：

\begin{enumerate}
\item \textbf{光谱分布的非均匀性得到初步实证}：仅19.8\%的美国企业使用了AI——这意味着超过80\%的企业停留在本文定义的0级（完全未使用）。而在使用AI的企业中，60\%的顶级采用率（信息技术行业）与5.3\%的最低采用率之间的差距，恰恰对应了光谱从1-2级到4-5级的分化。
\item \textbf{企业规模与光谱层级的正相关}：大企业35.3\% vs 小企业21.5\%的采用率差异，支持了本文关于''AI使用层级与社会既有分层高度相关''的判断。大企业拥有更多的资源进行AI培训、部署和定制——这对应于光谱爬升所需的外部推动力。
\item \textbf{生产力增益验证了光谱爬升的价值}：23-33\%的效率提升和每日40-60分钟的节省，说明从0级到2级以上的跃迁带来的经济效益是显著的。这为本文的政策建议——通过教育和开源降低光谱爬升门槛——提供了经济学的成本-收益论证。
\end{enumerate}

高盛在报告中写道：''我们继续观察到生成式AI在已部署的有限领域对劳动生产率产生巨大影响。''紧接着指出了核心矛盾——''正在使用AI的公司正在拉开差距，而他们的大多数竞争对手甚至还没有进入赛道。''这正是本文第2.6节所描述的光谱固化机制在产业层面的精确映射。

\subsubsection{中国官方数据：产业规模与普及率目标}

中国信息通信研究院2026年2月发布的《全球人工智能发展白皮书》提供了中国侧的宏观数据：2024年中国AI核心产业规模突破9000亿元，大模型综合性能显著提升，应用门槛与使用成本持续降低。白皮书将开源列为全球AI发展的关键路径，并指出中国开源大模型数量位居全球第二。

更具战略意义的是国务院2025年8月发布的《关于深入实施''人工智能+''行动的意见》。该文件设定了明确的量化目标：
\begin{itemize}
\item \textbf{到2027年}：AI与6大重点领域广泛深度融合，新一代智能终端、智能体应用普及率超过70\%
\item \textbf{到2030年}：AI全面赋能高质量发展
\end{itemize}

70\%的普及率目标——如果实现——意味着在政策推动下，中国社会的AI认知光谱将从当前的金字塔形（多数在0-2级）向更均衡的分布快速演化。这是本文定义的''光谱扩散优化''（情景B）在国家级政策层面的具体体现。需要指出，政策目标不等于现实——但政策目标塑造了资源分配的方向，而资源分配的方向影响着光谱演化的轨迹。

\subsection{综合：理论框架的多源实证支撑}

将以上智库和金融机构的最新报告与本文的理论框架进行对照，可以形成一张多源验证矩阵：

\begin{center}
\begin{tabular}{c|c|c}
\hline
本文核心论点 & 验证来源 & 验证类型 \\
\hline
AI使用光谱分布非均匀 & Goldman 19.8\%采用率数据 & 定量实证 \\
开源降低光谱爬升门槛 & CNBC: 中国模型便宜60-90\%，美企token占比飙至46\% & 市场行为验证 \\
闭源形成守门人结构 & RAND: 开源作为软实力；OpenAI应政府要求限发 & 政策分析验证 \\
光谱爬升需要外部推动 & Goldman: 大企业35.3\% vs 小企业21.5\% & 结构性验证 \\
认知民主化是国际竞争维度 & Brookings: 多维AI竞赛；国务院: 70\%普及率目标 & 战略分析验证 \\
开源的地缘技术意义 & RAND: 美国若限制开源可能将生态让给中国 & 地缘政治验证 \\
\hline
\end{tabular}
\end{center}

这一多源验证矩阵的意义在于：本文的核心论证在定量实证（Goldman数据）、市场行为（CNBC报道）、政策分析（Brookings/RAND）和战略规划（中国国务院目标）四个维度上均获得了独立支撑。这不仅增强了论证的可信度，也表明本文的理论框架具有跨越微观（企业采用率）、中观（产业竞争）和宏观（国家战略）的分析适用性。
